% Generated draft: requires mathematical and source review.
\documentclass[11pt]{article}
\usepackage[margin=25mm]{geometry}
\usepackage{fontspec}
\setmainfont{Latin Modern Roman}
\usepackage{amsmath,amssymb,mathtools,unicode-math}
\setmathfont{Latin Modern Math}
\usepackage{enumitem,needspace,xurl,hyperref}
\setlength{\parindent}{0pt}
\setlength{\parskip}{5pt}
\setlength{\emergencystretch}{3em}
\newcommand{\N}{\mathbb{N}}
\newcommand{\Z}{\mathbb{Z}}
\newcommand{\Q}{\mathbb{Q}}
\newcommand{\R}{\mathbb{R}}
\newcommand{\C}{\mathbb{C}}
\newcommand{\F}{\mathbb{F}}
\newcommand{\A}{\mathbb{A}}
\newcommand{\PP}{\mathbb{P}}
\newcommand{\E}{\mathbb{E}}
\newcommand{\eps}{\varepsilon}
\newcommand{\dd}{\,\mathrm{d}}
\begin{document}
\section*{UnsolvedMath ID: 30002715}
\textbf{Source problem code:} catalogue entry
\subsection{Definitions and mathematical statement}
Let \(K\) be a finitely generated field, understood in the usual sense as finitely generated over its prime field. The supplied question concerns every pair of \(K\)-varieties \(X,Y\) with \(\dim X\geq 1\) and \(\dim Y\geq 1\). To interpret their function fields, suppose that ``variety'' means an integral separated scheme of finite type over \(K\); the record does not explicitly specify this convention or impose geometric integrality, smoothness, or normality. Write \(L_X=K(X)\) and \(L_Y=K(Y)\).

Define \(\operatorname{Isom}_K(L_X,L_Y)\) to be the set of field isomorphisms \(L_X\to L_Y\) fixing \(K\). Let \(\operatorname{Bir}_K(Y,X)\) be the set of birational maps \(Y\dashrightarrow X\) defined over \(K\), with rational maps identified when they agree on a common dense open subset. Pullback reverses direction: a birational map \(f:Y\dashrightarrow X\) induces \(f^*:L_X\to L_Y\). This reversal must be accounted for when interpreting the source's first canonical map.

For a field \(L\), choose a separable closure \(L^{\mathrm{sep}}\) and write
\[
G_L=\operatorname{Gal}(L^{\mathrm{sep}}/L),
\]
with its profinite topology. One possible meaning of an outer Galois-group isomorphism is an element of
\[
\operatorname{OutIsom}(G_{L_Y},G_{L_X})
=
\operatorname{Isom}_{\mathrm{cont}}(G_{L_Y},G_{L_X})
\big/\operatorname{Inn}(G_{L_X}),
\]
where the quotient identifies isomorphisms differing by conjugation in the target. A field isomorphism \(L_X\to L_Y\) induces such an outer isomorphism in the displayed reverse direction.

\textbf{Question recorded by the source.} For every \(K,X,Y\) as above, are the canonical maps from \(K\)-isomorphisms of function fields to \(K\)-birational isomorphisms and then to the corresponding outer Galois-group isomorphisms bijective? With directions made explicit, the requested pattern is
\[
\operatorname{Isom}_K(L_X,L_Y)
\longrightarrow
\operatorname{Bir}_K(Y,X)
\longrightarrow
\mathcal{O}_K(G_{L_Y},G_{L_X}),
\]
where \(\mathcal{O}_K\) denotes the intended target in the source, whose definition is unresolved. In particular, the supplied record does not say whether this target is the full outer-isomorphism set defined above or a relative version requiring compatibility with restriction to \(G_K\). If compatibility is intended, the required equality or equality up to conjugacy, and the permitted inner conjugations, must be specified. Thus the display records the requested assertion but does not supply a fully determined final target.

The record also leaves positive-characteristic and inseparability qualifications unspecified: it does not say whether ordinary function fields or their perfect closures are intended, or whether additional restrictions on \(K\), \(X\), and \(Y\) are required. These omissions prevent a unique precise formulation of the full conjecture. The background additionally mentions broader anabelian-variety and generic-point assertions, but supplies no mathematical formulations of them; they are not additional statements that can be formalized from this record.
\subsection{Short English statement}
The source asks whether function-field isomorphisms, birational maps, and suitable outer Galois-group isomorphisms determine one another for positive-dimensional varieties over finitely generated fields. A precise final conjecture requires clarification of the Galois-group target, compatibility with the base field, and inseparability conventions. The supplied \texttt{partially\_solved} label is dated curation, not an independently verified conclusion.
\subsection{Sources}
\begin{itemize}
\item Local UnsolvedMath registry: \texttt{data/unsolvedmath/problems.json}, permanent ID \(30002715\), title ``Birational Anabelian Reconstruction over Finitely Generated Fields,'' source problem code \texttt{catalogue entry}. No individual proposer is supplied; the proposed year is recorded as \(2014\).
\item Supplied report attribution: \emph{Valuation Theory and Its Applications (2015)}, Oberwolfach Reports, pp.~2757--2823, DOI \texttt{10.4171/owr/2014/49}. Supplied source and literature link: \url{https://doi.org/10.4171/owr/2014/49}. The record attributes the problem to workshop open-problem discussions. The citation's \(2015\) year and its \(2014\) DOI component and proposed year are retained without resolving their relationship.
\item Supplied literature reference, labeled ``Pop, birational anabelian geometry'': \url{https://annals.math.princeton.edu/1994/139-1/p05}. The record supplies no fuller bibliographic details or theorem hypotheses.
\item Supplied literature assessment dated exactly \texttt{2026-08-21}: status \texttt{partially\_solved}. The assessment says that Pop proved the birational anabelian conjecture over finitely generated fields, while the broader anabelian-variety and generic-point assertions mentioned in the record are not all known in that generality. This is supplied dated curation derived from the Oberwolfach material, not independently verified here; the source itself requests independent verification of time-sensitive status and cited claims.
\item The registry labels the statement \texttt{corrected\_verified} and reports that it was checked against the official Oberwolfach report. This is a supplied verification claim, not a claim that the present draft has visited or checked either external source.
\end{itemize}
\end{document}
